\item
\uline{For} @@M0@@, @@M1@@ \uline{and} @@M2@@ \uline{are flat over} @@M3@@.
\item \uline{The sequence}
\sgaStableTarget{sga6.e12.eq.unnamed.occ2}@@M4@@
\uline{is exact and} @@M5@@ \uline{is an} @@M6@@\uline{-module flat over} @@M7@@.
\end{enumerate}
\uline{Then} @@M8@@ \uline{is representable by a quasi-affine morphism of finite presentation.}

Indeed, first note that 1)-4) are stable under every base change @@M9@@. Only the stability of 3) and 4) is not evident. But, by the Künneth formula (EGA III 6.9.9), 1) and 3) imply that the formation of the @@M10@@ and @@M11@@ commutes with base changes @@M12@@, which immediately implies the stability of conditions 3). Moreover, the formation of the cokernel of @@M13@@ therefore also commutes with base changes; since it is flat over @@M14@@ (by 4)), @@M15@@ is flat over @@M16@@ and commutes with base changes. Consequently, @@M17@@ commutes with base changes @@M18@@ and is therefore universally equal to @@M19@@.

The stability of the hypotheses implies two things: a. Since the assertion to be proved is local for f.p.p.f., by descent we may suppose that @@M20@@ has a section @@M21@@, so @@M22@@ has the section @@M23@@. b. To establish the assertion it suffices to consider an @@M24@@-point of @@M25@@, that is, by a) and 2), an invertible sheaf @@M26@@ on @@M27@@ equipped with a rigidification @@M28@@ relative to @@M29@@, and to show that the f.p.p.f. sheaf @@M30@@, the inverse image of this point in @@M31@@, is representable by a quasi-affine @@M32@@-scheme of finite presentation over @@M33@@.

